% !PRINTED 72
\weberpage{72}
%<<CONT
und also wird nach 3. ein Theiler $\Omega$ von $\Omega_m$ dann und nur
dann nicht primär sein, wenn in der zu $\Omega$ gehörigen Gruppe $\mathfrak{A}$
eine der Gruppen $Q_1, Q_2, \ldots Q_\tau$ enthalten ist. Wenn wir also
solche Theiler der Gruppe $\mathfrak{N}$, die keine der Gruppen $Q_1, Q_2 \ldots Q_\tau$
als Theiler enthalten, primäre Theiler nennen, so können wir
den Satz aussprechen:

\gesperrt{5. Um alle primären Theiler von $\Omega_m$ zu erhalten,
hat man alle primären Theiler der Gruppe $\mathfrak{N}$
aufzusuchen und die zugehörigen Körper zu
bilden.}

\gesperrt{6. Wenn man aber alle primären Theiler aller
vollen Kreistheilungskörper aufstellt, so erhält
man jeden Kreistheilungskörper, und jeden
nur einmal, und zwar jeden dargestellt durch
Einheitswurzeln möglichst niedrigen Grades.}

Der Körper $\Omega_{2m}$ hat, wenn $m$ ungerade ist, gar keinen primären Theiler und ist mit $\Omega_m$ identisch. In allen anderen Fällen
ist $\Omega_m$ wenigstens sein eigener primärer Theiler.

Es ist für das Folgende noch erforderlich, dass wir uns
über den Umfang und die Constitution der Gruppen $Q$ klar
werden.

\begin{center}
\S.\,19.

\gesperrt{Primäre und nicht primäre Theiler der Gruppe $\mathfrak{N}$.}
\end{center}

Wir bezeichnen mit $q$ irgend eine der in $m$ aufgehenden
Primzahlen und setzen $m = q m'$. Dann betrachten wir die
Gruppe
\[
Q \equiv 1\,(\mathrm{mod}\ m').
\]

Um die Bedingungen für eine Zahl in $Q$ zu ermitteln, wenden
wir die Darstellung der Gruppe $\mathfrak{N}$ durch eine Basis und die
Bezeichnungsweise der Indices an, wie wir sie im §. 16 eingeführt
und erklärt haben, und bezeichnen danach die zu $Q$ reciproke
Gruppe mit $P$.

Die Indices einer Zahl in $Q$ bezeichnen wir mit
\[
\gamma_{-1}, \gamma_0, \gamma_1, \ldots \gamma_\mu,
\]
und einer Zahl in $P$ mit
\[
\delta_{-1}, \delta_0, \delta_1, \ldots \delta_\mu.
\]
